% Draft for sec6_children.tex, paragraph 'What the model takes from this' (WP7 of fable_10h_plan_paperA.md, 2026-10-20).
% NOT input by main.tex; for the lead to merge. Numbers: IGE/output/parents_present_outcomes.md (code
% IGE/79.ParentsPresentOutcomes.py). Bib keys used exist in paperA_joint_loss/references.bib.
%
% Hilger numbers. The AEJ:Applied article (8(3), 247-283) is behind the AEA paywall and was not read. The abstract
% says layoffs "dramatically reduce family income but only slightly reduce college enrollment, college quality, and
% early career earnings". Numbers from Hilger's 2013 Harvard dissertation, chapter 1 (the same design; DASH
% urn-3:HUL.InstRepos:11129107), which is the paper's earlier version:
%   - short-run post-tax family income loss about $8,000 (a 13.7% loss, as used by the CES-WP-24-15 calibration);
%     permanent family income about -$6,000;
%   - college enrollment at 18-22 -0.43 pp (s.e. 0.11); total college years at 18-22 -0.016;
%   - alumni-earnings-based college quality -0.35%; implied lifetime earnings loss $800-1,600 (fn. 34), i.e.
%     about 0.1-0.3% of lifetime earnings.
% Huttunen and Riukula (IZA DP 12788) summarise the AEJ result as "a very small effect ... on educational attainment
% and no effect on earnings". The low slope below uses -0.35% (the larger of the two dissertation numbers).
% TODO before coding: check it against the AEJ earnings table, as the identification plan requires.

\paragraph{What the model takes from this.} The PSID associations are an upper bound on the total effect of a father's
job loss on his children. They are robust only for White children's schooling and for job losses during childhood.
Administrative data find smaller effects: about $-9\%$ in adult earnings in Canada
\citep{oreopoulos2008intergenerational}, and much less in US tax data. There, a father's layoff during adolescence
cuts family income sharply, lowers college enrollment by less than half a percentage point, and leaves early-career
earnings almost unchanged \citep{hilger2016parental}. The model's child block maps family income at ages 0--5 and
6--17, and the share of childhood spent with both parents, into adult log earnings. The income slopes are set so that
a father's displacement at the mean displacement age lowers his child's adult earnings by 9\% (the high slope) or by
0.35\% (the low slope, from the US estimates). Counterfactuals are reported under both. The presence slope comes from
the PSID. Holding family income at both childhood ages fixed, a childhood spent entirely with both parents is not
associated with higher adult earnings ($+0.011$, s.e.\ 0.045, $p=0.80$). It is associated with 0.30 more years of
schooling ($p<0.001$). Holding fixed only the family's income in the three years before and after the father's
displacement, the association with adult earnings is $0.19^{***}$. The model uses 0.19 as the high presence slope
and zero as the low one.\footnote{PSID children born 1968--1995 (SRC and SEO, unweighted) with at least three
childhood interviews. ``Both parents present'' means the head and the wife of the child's family unit are the child's
mother and father at birth. The share of childhood with both parents averages 0.38 for Black and 0.70 for White
children. Adult earnings are the mean over the years as head or wife at 25--55 (at least three years; zeros kept).
All regressions control for sex, birth year, parental education and, in pooled rows, race. Standard errors are
clustered by 1968 family; siblings are not compared. Without income controls, a childhood with both parents is
associated with adult earnings 29.5 log points higher ($0.295^{***}$, s.e.\ 0.041; $N=5{,}091$), with an earnings
rank 9.5 percentiles higher ($9.46^{***}$) and with 0.72 more years of schooling ($0.724^{***}$). Net of log family
income at 0--5 and 6--17, the earnings coefficient is $0.011$ (s.e.\ 0.045, $p=0.80$; $N=4{,}701$): Black $-0.031$
($p=0.65$), White $0.069$ ($p=0.24$). The schooling coefficient is $0.299^{***}$ (s.e.\ 0.081; $N=6{,}213$): Black
$0.213^{*}$ ($p=0.054$), White $0.476^{***}$. The income slopes in that regression are $0.148^{***}$ (ages 0--5) and
$0.320^{***}$ (6--17) for log earnings; they are associations, which is why the model takes its income slopes from
administrative estimates instead. Splitting the shares by age, the share with both parents at 6--17 is positive for earnings
($0.102^{**}$, $p=0.038$) and schooling ($0.342^{***}$). The share at 0--5 is negative for earnings ($-0.126^{**}$,
$p=0.032$) and not significant for schooling ($-0.080$, $p=0.41$). The two age shares are strongly collinear, and
their sum is close to the undivided coefficient. In the father's
displacement sample of Section~\ref{sec:children} (pooled, $N=3{,}520$), the regression of adult log earnings on the
displacement, the share with both parents, and log family income before and after the (pseudo) event gives
displacement $-0.040$ ($p=0.25$), presence $0.187^{***}$ (s.e.\ 0.054), income before $0.179^{***}$ and income after
$0.209^{***}$. Black and White estimates are $0.178^{**}$ ($p=0.037$) and $0.175^{**}$ ($p=0.013$). Without the
income after the event, the presence coefficient is $0.242^{***}$ and the displacement coefficient $-0.084^{**}$
($p=0.016$). Code: \texttt{IGE/79.ParentsPresentOutcomes.py}.} The parents' side of the model is estimated without
the child block, so its results for consumption and for the ranking of policies stand on their own.

% ---------------------------------------------------------------------------------------------------------------
% Replacement text for IGE/proposals/paperA_identification_plan.md line 259 ('Child-block slopes' row), not edited:
%   high: beta_05 = beta_617 = -0.09 / Delta ln y_bar(sim) (Oreopoulos et al. 2008; Page et al. 2009); low: the same
%   with -0.0035 in place of -0.09 (Hilger: alumni-earnings college quality -0.35% for a layoff that cuts family
%   income ~13.7%; check against the AEJ earnings table); beta_P: high 0.19 = PSID presence slope net of family
%   income before and after the father's displacement (IGE/79 spec E2, pooled, s.e. 0.054), low 0 = net of family
%   income at 0-5 and 6-17 (IGE/79 spec B, pooled 0.011, s.e. 0.045, p=0.80). Schooling analogue 0.30*** (spec B).
%   Script: IGE/79.ParentsPresentOutcomes.py.
% Text for counterfactuals.py docstring 'Children's adult outcomes' (IGE/paperA_model is frozen until the CMA-ES
% stage): beta_05 = beta_617 = {-0.09 (high), -0.0035 (low)} / Delta ln y_bar(sim); beta_P = {0.19 (high), 0 (low)}
% per unit share of childhood with both parents (IGE/79).
